\section{Example}

Given the following equation:
\[A B (\vec{v} + \vec{w})\ \vert\ v, w \in \R^3, A, B \in \R^{3\times3}\]

The code below recreates the equation with exemplary values and directly calculates the results.
Then, the results are used to print a step-by-step solution with the chosen values.\\
Note:
\begin{itemize}
\setlength\itemsep{0em}
\item While creating a new vector or matrix is done by providing its name as a simple string, accessing its content after creation is done as a command (with backslash)
\item This package mainly provides macros for mathematical calculations. These macros assign the calculated result automatically to a variable using the optional parameter (for example \textbackslash addVec). If you wish to specify the name of the variable yourself, you may pass a string for the optional argument.
\end{itemize}

\lstset{style=latexStyle}
\begin{lstlisting}[gobble=4]
% Initialize two vectors and add them together
\newVec{myVecV}{1,2,3}
\newVec{myVecW}{7,8,9}
\addVec[addVecResult]{\myVecV}{\myVecW}

% Initialize two matrices and multiply them with each other
\newMat{myMatA}{{{0.5,2,3}{4,5,6}{7,8,9}}}
\newDiagMat{myMatB}{2, 3, 4}
\mulMat[mulMatResult]{\myMatA}{\myMatB}

% Transform the result vector with the created matrix
\transformVec{\mulMatResult}{\addVecResult}

% Print a step-by-step solution
\begin{align}
    \printMat{\myMatA} \printMat{\myMatB}
    \left(\printVec{\myVecV} + \printVec{\myVecW}\right)
    =
    \printMat{\mulMatResult}
    \printVec{\addVecResult}
    =
    \printMat{\transformVecResult}
\end{align}
\end{lstlisting}
\vspace{5mm}

When compiling the above \LaTeX{} code, the following output is created.
\begin{align*}
    \begin{pmatrix}0.5&2&3\\4&5&6\\7&8&9\end{pmatrix}
    \begin{pmatrix}2&0&0\\0&3&0\\0&0&4\end{pmatrix}
    \left(
        \begin{pmatrix}1\\2\\3\end{pmatrix} +
        \begin{pmatrix}7\\8\\9\end{pmatrix}
    \right)
    = \begin{pmatrix}1&6&12\\8&15&24\\14&24&36\end{pmatrix}
    \begin{pmatrix}8\\10\\12\end{pmatrix}
    = \begin{pmatrix}212\\502\\784\end{pmatrix}
\end{align*}

All three classes provide various macros to print the content of a variable for all kinds of use cases.
For example, you can print a vector inside a line of text using the \textbackslash printVecT, which prints a transposed version of the vector's contents.

\clearpage